
    \begin{table}[h!]
    \centering
    \resizebox{\textwidth}{!}{%
    \begin{tabular}{p{0.2\textwidth} p{0.4\textwidth} p{0.4\textwidth} p{0.4\textwidth}}
    \toprule
    \textbf{Method Name} & \textbf{Metric Definitions} & \textbf{Compute Attribution} & \textbf{Reporting Standards} \\
\midrule
FastAdaSP\cite{lu2024fastadaspmultitaskadaptedefficientinfere} & Real Time Factor & Constant/decay Scheduling & A100 80gb \\
StreamUni\cite{guo2025streamuniachievingstreamingspeechtransla} & Al, Laal, Streamlaal & Delay K & Al, Sacrebleu, Comet \\
LSLM\cite{ma2024languagemodelcanlistenwhilespeaking} & Success Window 1S & Middle Fusion & Wer And F1 \\
SyllableLM\cite{baade2024syllablelmlearningcoarsesemanticunitsspe} & Rtf Measured & Unit Rates Controlled & Nvidia A40 Hardware \\

    \bottomrule
    \end{tabular}%
    }
    \caption{Comparison of representative recent systems across the three pillars of latency evaluation: metric definitions, compute attribution and scheduling, and reporting standards/hardware. For each method, the table records the primary latency/quality metrics (e.g., RTF, AL/LAAL/StreamLAAL, success-window), the scheduling or attribution policy (e.g., delay-k, mid-fusion, unit-rate/decay), and the reported baselines or devices (e.g., BLEU/COMET/WER, A100/A40). The heterogeneity highlights the need for the standardized measurement, scheduling, and reporting practices articulated in the accompanying framework.}
    
\label{tab:Arbitrary_table_1}\end{table}